\documentclass[../comorg-part1.tex]{subfiles}

\begin{document}
\section{Data Organisation in Memory}

\subsection{Memory Hierarchy}
\begin{figure}[H]
    \centering
    \includegraphics[width=0.8\textwidth]{images/memory-hierarchy.png}
    % \caption*{}
\end{figure}
\begin{enumerate}
    \item \textbf{Registers} Very fast access, but limited numbers within the CPU (2-128 registers). Operates at CPU clock rate.
    \item \textbf{Cache memory} Fast access static RAM close to CPU. Typical access time 3-20ns, size up to 512kB.
    \item \textbf{Main memory} Dynamic RAM or ROM. Typical access time 30-70nS, size up to 16GB.
    \item \textbf{Secondary memory} Not always random access but non-volatile. Maybe be based on magnetic or flash technology. Typical access time 0.03-100mS, size: up to 4TB.
\end{enumerate}

\subsubsection{Main Memory}
\begin{figure}[H]
    \centering
    \includegraphics[width=0.3\textwidth]{images/main-mem.png}
    \caption*{Main memory}
\end{figure}
\begin{itemize}
    \item Each address location is byte-size (8 bits).
    \item The memory size (the number of addresses) i.e. $2^n$ bytes, is dependent on the number of lines $n$ in the address bus.
    \item Based on Von Neumann's architecture, the memory stores both code (instructions) and data.
\end{itemize}

\subsection{Number representation}
\begin{table}[h!]
    \centering
    \begin{tabular}{r c c l} 
        \hline
        \textbf{Type} & \textbf{Bytes} & \textbf{Bits} & \textbf{\quad Range} \\ 
        \hline
        signed char & 1 & 8 & $-128 \to +127$ \\ 
        unsigned char & 1 & 8 & $0 \to +255$ \\ 
        short int & 2 & 16 & $-32,768 \to +32,767 \quad (\pm 32\text{KB})$ \\ 
        unsigned short int & 2 & 16 & $0 \to +65,535 \quad\qquad (64\text{KB})$ \\ 
        unsigned int & 4 & 32 & $0 \to +4,294,967,295 \quad (4\text{GB})$ \\ 
        int & 4 & 32 & $-2,147,483,648 \to +2,147,483,647 \ (\pm 2\text{GB})$ \\ 
        long int & 4 & 32 & $-2,147,483,648 \to +2,147,483,647 \ (\pm 2\text{GB})$ \\ 
        long long int & 8 & 64 & $-(2^{63}) \to (2^{63})-1$ \\ 
        float & 4 & 32 &  \\ 
        double & 8 & 64 &  \\ 
        long double & 12 & 96 &  \\ 
        \hline
    \end{tabular}
\end{table}
Note:
\begin{itemize}
    \item Integers can be signed or unsigned.
    \item Floating point numbers are always signed.
\end{itemize}

\subsubsection{Two ways of storing a multi-byte number}
\begin{figure}[H]
    \centering
    \includegraphics[width=0.9\textwidth]{images/32bit-Endianess.png}
    \caption*{Little vs Big Endian format}
\end{figure}
\begin{itemize}
    \item Little Endian - LSB at lowest address.
    \item Big Endian - MSB at lowest address.
\end{itemize}

\subsection{Character Representation}
A char requires one byte (8 bits) of memory.
The binary data is transformed to some representative characters through some encoding standard.\newline
Some encoding standards:
\begin{itemize}
    \item 7-bit ASCII code
    \item DEC's Sixbit (6-bit)
    \item IBM's EBCDIC (8-bit)
    \item Unicode (16-bit)
\end{itemize}
\begin{figure}[H]
    \centering
    \includegraphics[width=0.9\textwidth]{images/ascii.png}
    \caption*{7-bit ASCII table}
\end{figure}
Note that lower and upper case characters and digits are contiguous.\newline

\bigskip

\noindent Boolean variables have 2 states:
\begin{enumerate}
    \item False = 0
    \item True = non-zero
\end{enumerate}
\noindent Most implementations use one byte (8 bits) to store a 1-bit Boolean value (inefficient).
Some processors like the 8051 support a small portion of bit-addressable memory.
\begin{figure}[H]
    \centering
    \includegraphics[width=0.4\textwidth]{images/8051.png}
    \caption*{8051 processor supports bit-addressable memory.}
\end{figure}

\subsection{Array, String and Structure Representations}
\subsubsection{Arrays}
A linear array is a consecutive area in memory storing a homogenous data type.\newline
Elements are accessed through some offset from the base address.
\begin{figure}[H]
    \centering
    \includegraphics[width=0.9\textwidth]{images/array-storage.png}
    \caption*{Offset calculated from number of elements and size of the data type.}
\end{figure}

\begin{figure}[H]
    \centering
    \includegraphics[width=0.9\textwidth]{images/nested-array.png}
    \caption*{Nested arrays}
\end{figure}

\subsubsection{Strings}

\begin{figure}[H]
    \centering
    \includegraphics[width=0.5\textwidth]{images/c-string.png}
    \caption*{Strings are initialised as array of characters.}
\end{figure}


\begin{figure}[H]
    \centering
    \includegraphics[width=0.6\textwidth]{images/c-vs-pascal-string.png}
    \caption*{C string terminates with null character, 
    Pascal string stores the length of the string at the start of the string.}
\end{figure}
\begin{itemize}
    \item C strings have no length limit, whereas Pascal strings are limited to number of bytes used to store the length (e.g. 1 byte $\rightarrow$ max 255 characters.)
    \item C strings have no overhead (only one byte for \textbackslash 0) for length metadata, Pascal strings require space to store the length prefix.
    \item Determining length of C string is O(n), for Pascal string is O(1).
\end{itemize}

\subsubsection{Structures}

\begin{figure}[H]
    \centering
    \includegraphics[width=0.9\textwidth]{images/struct-representation.png}
    \caption*{Each data type is ordered according to declaration, and allocated their own data size based on their type.}
\end{figure}

\subsection{Pointer representation}

\begin{figure}[H]
     \centering
     \begin{subfigure}[b]{0.45\textwidth}
         \centering
         \includegraphics[width=\textwidth]{images/ptr-code.png}
         \caption*{Initialising pointer}
     \end{subfigure}
     \hfill
     \begin{subfigure}[b]{0.45\textwidth}
         \centering
         \includegraphics[width=\textwidth]{images/ptr-visual.png}
         \caption*{Assuming addresses are 16 bits.}
     \end{subfigure}
        \caption{Size of pointer is fixed regardless of data type.}
\end{figure}
\noindent Use \texttt{*ptr} to dereference a pointer.

\subsection{Data Alignment}
\begin{table}[H]
    \centering
    \renewcommand{\arraystretch}{1.5} % Increases row height for readability
    \begin{tabular}{|l|c|l|}
        \hline
        % \rowcolor[HTML]{B4D5E1} % Approximate light blue from your slide
        \textbf{Data Type} & \textbf{Size (Byte)} & \textbf{Example of allowable start addresses due to alignment} \\ \hline
        char    & 1 & 0x..0000, 0x..0001, 0x..0002 \\ \hline
        short   & 2 & 0x..0000, 0x..0002, 0x..0004 \\ \hline
        int     & 4 & 0x..0000, 0x..0004, 0x..0008 \\ \hline
        float   & 4 & 0x..0000, 0x..0004, 0x..0008 \\ \hline
        double  & 8 & 0x..0000, 0x..0008, 0x..0010 \\ \hline
        pointer & 8 & 0x..0000, 0x..0008, 0x..0010 \\ \hline
    \end{tabular}
    \caption*{Multi-byte data should be aligned to addresses that are multiples of their data type size.}
\end{table}
\textit{Assuming} a 64-bit processor that uses 64-bits to represent an address.

\begin{figure}[H]
    \centering
    \includegraphics[width=0.8\textwidth]{images/intel-pentium-4.png}
    \caption*{Intel Pentium 4}
\end{figure}
\begin{itemize}
    \item 64-bit data bus $\rightarrow$ The CPU fetches 64-bits every memory cycle.
    \item 33 address pins $\rightarrow$ There are $2^33$ 8 byte "blocks".
    \item 8 Bank Enable pins select which 8-bits in that 64-bits block to read from.
    \item The CPU has a 64-bit data word size.
\end{itemize}
\begin{figure}[H]
    \centering
    \includegraphics[width=0.3\textwidth]{images/8-byte-alignment.png}
    \caption*{8-byte data should be aligned to addresses that are multiples of 8.}
\end{figure}
\begin{itemize}
    \item Since the CPU fetches 8 bytes at a time, any 8-byte data that is not aligned
    needs two memory cycles to be fetched.
    \item If a processor has an 8 bit (1 byte) data bus, there won't be any alignment issues.\newline
    Multi-byte data takes the same number of cycles to be fetched, regardless of the memory address.
    \item If a processor has a 16 bit (2 bytes) data bus, then 2, 4 or 8 byte data
    should be aligned to even addresses.
\end{itemize}

\subsubsection{Data alignment in structures}
\begin{figure}[H]
     \centering
     \begin{subfigure}[b]{0.25\textwidth}
         \centering
         \includegraphics[width=\textwidth]{images/struct-1.png}
         \caption*{Defining struct in C}
     \end{subfigure}
     \hspace{1cm}
     \begin{subfigure}[b]{0.65\textwidth}
         \centering
         \includegraphics[width=\textwidth]{images/struct-padding.png}
         \caption*{Padded bytes are added to maintain alignment for each data type.}
     \end{subfigure}
        % \caption{}
\end{figure}

\begin{figure}[H]
     \centering
     \begin{subfigure}[b]{0.3\textwidth}
         \centering
         \includegraphics[width=\textwidth]{images/struct-2.png}
         \caption*{Rearranging order of data objects}
     \end{subfigure}
     \hspace{1cm}
     \begin{subfigure}[b]{0.3\textwidth}
         \centering
         \includegraphics[width=\textwidth]{images/struct-optimised.png}
         \caption*{New struct takes 8 bytes instead of 12 previously.}
     \end{subfigure}
        % \caption{Size of pointer is fixed regardless of data type.}
\end{figure}



\end{document}